\documentclass[11pt, a4paper]{article}

\usepackage[T1]{fontenc}
\usepackage[utf8]{inputenc}
\usepackage{tgpagella}
\usepackage[spanish, es-noshorthands]{babel}
\usepackage{amsmath, amssymb, bm}
\usepackage{tabularx}
\usepackage{booktabs}
\usepackage[most]{tcolorbox}
\usepackage[margin=2.5cm]{geometry}
\usepackage{ragged2e}
\usepackage{fancyhdr}

\pagestyle{fancy}
\fancyhf{}
\setlength{\headheight}{16pt}
\lhead{\textbf{UCB Sede Tarija} | Ing. Mecatrónica}
\chead{}
\rhead{IMT-121 Circuitos Electrónicos I | Checklist Final Lab}
\lfoot{Prof: M.Sc. Ing. Bernardo Quiroga Turdera}
\rfoot{Página \thepage}
\renewcommand{\headrulewidth}{0.4pt}

\definecolor{ucbblue}{RGB}{0, 51, 102}

\begin{document}

\begin{tcolorbox}[colback=ucbblue!5!white, colframe=ucbblue, title=\textbf{Checklist de Integración: Cierre Total de Laboratorio (A+B)}]
    \centering \textbf{Preparación de Evidencia para Reporte IEEE}
\end{tcolorbox}

\noindent \textit{Instrucción: Cada equipo debe asegurar que la evidencia primaria listada a continuación esté completamente registrada en sus bitácoras u hojas de datos ANTES de comenzar la redacción del reporte IEEE. Un faltante impedirá completar técnicamente el informe.}

\vspace{0.3cm}
\textbf{Equipo:} \hrulefill \quad \textbf{Validado por:} \hrulefill

\section*{1. Experiencia Integrada A}
\begin{itemize}\setlength\itemsep{0.1em}
    \item[$\square$] Evidencia de \textbf{Superposición} completa (Teoría, SPICE, Banco).
    \item[$\square$] Medición física explícita de $V_{oc}$.
    \item[$\square$] Determinación segura (red desenergizada) de $R_{th}$.
    \item[$\square$] Parámetros completos de equivalentes Thévenin y Norton ($I_N$).
    \item[$\square$] Comparación directa SPICE - Medición con cálculo de Error Relativo.
    \item[$\square$] Identificación y documentación de \textbf{al menos una discrepancia técnica}.
\end{itemize}

\section*{2. Experiencia Integrada B}
\begin{itemize}\setlength\itemsep{0.1em}
    \item[$\square$] Barrido paramétrico sobre múltiples valores de $R_L$.
    \item[$\square$] Registro de $V_L$ e $I_L$ para cada caso en la tabla.
    \item[$\square$] Cálculo consistente de $P_L$.
    \item[$\square$] Trazo o representación gráfica de la \textbf{Curva de Potencia} (campana).
    \item[$\square$] Identificación explícita de la resistencia $R_L$ donde ocurrió el \textbf{máximo numérico}.
    \item[$\square$] Comparación de dicho máximo experimental contra el valor de $R_{th}$.
    \item[$\square$] Análisis del Efecto de Carga (Loading Effect) documentado.
    \item[$\square$] Incorporación de la comparación Teoría-SPICE-Medición para el análisis de potencia.
\end{itemize}

\vspace{0.4cm}
\begin{tcolorbox}[colback=red!5!white, colframe=red!70!black, title=\textbf{Pregunta Final de Control de Calidad}]
¿Qué \textbf{único} elemento faltante, cálculo omitido o discrepancia no explicada en esta lista podría impedir que su reporte IEEE esté técnicamente completo y defendible?
\vspace{1cm}
\end{tcolorbox}

\end{document}
